\documentclass[11pt,a4paper]{article}
\usepackage[margin=1.1in]{geometry}
\usepackage{amsmath,amssymb,amsthm,mathtools}
\usepackage[T1]{fontenc}
\usepackage{microtype}
\usepackage{enumitem}
\usepackage{booktabs}
\usepackage[hidelinks]{hyperref}
\hypersetup{pdfauthor={Dritero Mehmetaj},pdftitle={A finite spectrum on an infinite lattice},pdfsubject={17B25}}

\newtheorem{theorem}{Theorem}[section]
\newtheorem{proposition}[theorem]{Proposition}
\newtheorem{lemma}[theorem]{Lemma}
\newtheorem{corollary}[theorem]{Corollary}
\theoremstyle{definition}
\newtheorem{definition}[theorem]{Definition}
\newtheorem{remark}[theorem]{Remark}
\theoremstyle{remark}
\newtheorem{scope}[theorem]{Scope}

\newcommand{\g}{\mathfrak{g}}
\newcommand{\esix}{\mathfrak{e}_6}
\newcommand{\slt}{\mathfrak{sl}_2}
\newcommand{\z}{\mathfrak{z}}
\newcommand{\Z}{\mathbb{Z}}
\newcommand{\Q}{\mathbb{Q}}
\newcommand{\C}{\mathbb{C}}
\newcommand{\F}{\mathbb{F}}
\newcommand{\ad}{\operatorname{ad}}
\newcommand{\im}{\operatorname{im}}
\newcommand{\charpoly}{\operatorname{charpoly}}
\newcommand{\SU}{\operatorname{SU}}

\title{A finite spectrum on an infinite lattice:\\
centraliser dimensions for a four-dimensional abelian\\
subalgebra of $\esix$, and why sampling cannot find them}
\author{Drit\"ero Mehmetaj}
\date{26 August 2026}

\begin{document}
\maketitle

\begin{abstract}
Let $C\subset\esix$ be the four-dimensional abelian subalgebra spanned by highest-weight
vectors for the principal $\slt$-summands of exponents $4,7,8,11$, with coefficients supplied
by the binary tetrahedral invariant forms of degrees $8,14,16,22$. We study the function
$S\mapsto\dim\z(S)$ on the subspace lattice of $C$.

That lattice is infinite, so no sample determines the function. We show its image is
nevertheless finite and exactly computable. The mechanism is a structure theorem, exact over
$\Q$: the adjoint action of $C$ makes $\esix=\z(C)\oplus V'$ with $C$ acting as \emph{literally
zero} --- not merely nilpotently --- on the $12$-dimensional summand $\z(C)$, and the $66$
non-zero weights fall into three Galois orbits, of sizes $6,12,12$ and multiplicities $3,1,3$.
Consequently, for \emph{every} subspace $S\subseteq C$,
\[
   \dim\z(S)\;=\;12+\sum\bigl\{\,m_\lambda \;:\; \lambda|_S=0\,\bigr\},
\]
so the function is the flat-function of an arrangement of $30$ hyperplanes in a $4$-dimensional
space. Enumerating the $109$ flats gives the image exactly:
\[
   \{12,\,14,\,16,\,18,\,20,\,26,\,28,\,30,\,36,\,46,\,78\},
\]
eleven values, every one attained. A prior sample --- sixteen coordinate subspaces together
with $440$ random rational directions --- had returned only $\{12,30,78\}$. We explain the
discrepancy: the eight missing values are supported on proper subvarieties, which a random
rational direction misses with probability one. The extreme case is $46$, which on a
distinguished plane occurs exactly on the roots of an irreducible rational cubic, and so is not
merely rare over $\Q$ but \emph{arithmetically inaccessible} over $\Q$; that cubic acquires a
root in, and generates, the cubic field $K=\Q[x]/(x^3-12x-5)$. The structure theorem, the
weight orbits, the master formula and the cubic are exact over $\Q$; the flat enumeration is
exhaustive at three faithful primes and is stated as such.
\end{abstract}

\begin{center}\small
\textit{Paper II of four.} The others are I, \emph{The period-one locus in $\mathrm{GL}(2,\mathbb{Z})$} · III, \emph{The boundary-graviton one-loop determinant} · IV, \emph{What a class invariant cannot supply}.
They share a verification convention: every numeral is recomputed by a named script whose
expected output is stated, and every control is designed so that a plausible error makes it fail.
\end{center}

\medskip\noindent\small\textbf{MSC 2020.} Primary 17B25; Secondary 52C35, 17B45, 20G41, 11R16.\\
\textbf{Keywords.} exceptional Lie algebras, centraliser, hyperplane arrangement, flat-function, binary tetrahedral group.\normalsize


\section{Introduction}

\subsection{The problem}

Let $\g$ be a semisimple Lie algebra and $C\subseteq\g$ an abelian subalgebra. For a subspace
$S\subseteq C$ write $\z(S)=\{X\in\g:[X,S]=0\}$. The assignment $S\mapsto\dim\z(S)$ is a
function on the subspace lattice of $C$, and when $\dim C\ge2$ that lattice is infinite: it
contains a projective space of lines. So the function cannot be tabulated, and any finite
computation is a sample.

Samples of this function are systematically misleading, and the reason is structural rather
than statistical. $\dim\z(S)$ is upper semicontinuous: it jumps \emph{up} on closed subvarieties
of the Grassmannian and takes its minimum on a dense open set. A finite sample of rational
points therefore sees the generic value, plus whichever special values happen to lie on
coordinate subspaces one thought to try. It cannot see the rest, and no enlargement of the
sample repairs this.

This paper computes the function exactly for one specific $C$, by replacing the sample with an
enumeration of a hyperplane arrangement.

\subsection{The subalgebra}\label{sec:C}

Under a principal $\slt\subseteq\esix$, the adjoint representation decomposes into irreducible
summands of dimensions $2m_i+1$ where $m_i\in\{1,4,5,7,8,11\}$ are the exponents of $E_6$
\cite{Kostant}; the dimensions are therefore $3,9,11,15,17,23$.

\begin{definition}\label{def:C}
Recall that the space of binary forms of degree $n$ in $x,y$ \emph{is} the irreducible
$\slt$-representation of dimension $n+1$, with $\slt$ acting by the usual differential
operators. Under the principal $\slt$, the summands of $\esix$ of exponents $4,7,8,11$ have
dimensions $9,15,17,23$, hence are the spaces of binary forms of degrees $8,14,16,22$
respectively.

Let
\[
   t=x^5y-xy^5\ (\deg 6),\qquad \Phi=x^8+14x^4y^4+y^8\ (\deg 8)
\]
be two of Klein's binary tetrahedral forms. The four products
\[
   \Phi\ (\deg 8),\qquad t\Phi\ (\deg 14),\qquad \Phi^2\ (\deg 16),\qquad t\Phi^2\ (\deg 22)
\]
have exactly the four degrees above, so each is canonically an element of the corresponding
$\esix$-summand. Write $x_8,x_{14},x_{16},x_{22}\in\esix$ for these elements and set
\[
   C=\langle x_8,\;x_{14},\;x_{16},\;x_{22}\rangle .
\]
No choice of normalisation is involved beyond that of the principal $\slt$ itself: the
identification of degree-$n$ forms with the $(n+1)$-dimensional summand is canonical up to the
overall scalar on each summand, and every statement below is invariant under rescaling the
individual $x_n$.
\end{definition}

The four degrees are not a selection. The binary tetrahedral group $2T\subset\SU(2)$ is the
McKay partner of $E_6$, and its invariant ring is generated in degrees $6,8,12$
\cite{Klein,McKay}. Comparing that list of realisable degrees with the degrees of the principal
$\slt$-summands of $\esix$ leaves no freedom:

\begin{proposition}[the subalgebra is forced by degree matching]\label{prop:forced}
Let $D=\{6a+8b+12c: a,b,c\ge0\}$ be the set of degrees of $2T$-invariant binary forms, and let
$E=\{2,8,10,14,16,22\}$ be the degrees of the principal $\slt$-summands of $\esix$. Then
\[
   D\cap E=\{8,\,14,\,16,\,22\},
\]
the excluded degrees being $2$ and $10$ --- the exponents $1$ and $5$ --- which are not degrees
of any $2T$-invariant. Moreover each element of $D\cap E$ is realised by a \emph{unique}
monomial in the generators, namely
\[
   8\mapsto\Phi,\qquad 14\mapsto t\Phi,\qquad 16\mapsto\Phi^2,\qquad 22\mapsto t\Phi^2 ,
\]
none of which involves the degree-$12$ generator. Hence $C$ of Definition~\ref{def:C} is, up to
scaling each summand, the \emph{unique} subspace of $\esix$ spanned by $2T$-invariant forms
sitting in principal $\slt$-summands.
\end{proposition}

\begin{proof}
$D$ contains every even integer $\ge6$ except $10$: $6,8,12,14,16,18,\dots$ are realised by
$t,\Phi,t^2,t\Phi,\Phi^2,t^3,\dots$, while $10=6a+8b+12c$ has no solution in non-negative
integers, and neither does $2$. Intersecting with $E$ gives $\{8,14,16,22\}$. Uniqueness of the
monomial: for $8$, $6a+8b+12c=8$ forces $(a,b,c)=(0,1,0)$; for $14$, $(1,1,0)$; for $16$,
$(0,2,0)$; for $22$, $(1,2,0)$ --- in each case by inspection, since $c\ge1$ already exceeds the
target in the first two cases and leaves $4$ or $10$, neither in $D\cup\{0\}$, in the others.
\end{proof}

\begin{remark}[the partner does the work]
Running the same degree-intersection for the other two binary polyhedral groups shows the
matching is not a coincidence of arithmetic that any group would produce. Against $\esix$'s
summand degrees $\{2,8,10,14,16,22\}$:
\[
   D(2T)\cap E=\{8,14,16,22\},\qquad D(2O)\cap E=\{8,16\},\qquad D(2I)\cap E=\varnothing .
\]
The $E_6$ partner yields a four-dimensional $C$; the $E_7$ partner would yield only a
two-dimensional one, and the $E_8$ partner nothing at all.
\end{remark}

\begin{remark}
So the construction reads: \emph{take the finite group whose McKay partner is $E_6$, and place
its invariants inside $\esix$ wherever the degrees allow}. That determines $C$ completely. This
does not make $C$ canonical in an absolute sense --- it still depends on the principal $\slt$,
which is an input --- but it removes the suspicion that four summands were picked to obtain a
convenient answer.
\end{remark}

\begin{proposition}\label{prop:abelian}
$C$ is a four-dimensional abelian subalgebra of $\esix$: all six pairwise brackets vanish
identically, and $\dim\z(C)=12$. Both statements are exact over $\Q$.
\end{proposition}

The exponents used are $4,7,8,11$ --- all six exponents of $E_6$ except $1$ and $5$, these
being exactly the ones for which the tetrahedral products above have matching degree.
Nothing below depends on \emph{why} this particular $C$ is of interest; the results are
statements about the $C$ of Definition~\ref{def:C}, and the reader may take it as given.

\subsection{Results}

\begin{itemize}[leftmargin=1.4em]
\item \textbf{Theorem~\ref{thm:structure}} (structure, exact over $\Q$). $\esix=\z(C)\oplus V'$
      with $C$ acting as zero on $\z(C)$; the characteristic polynomial of a generic
      $R\in\ad(C)$ is $t^{12}q_6^{3}q_{12}q_{12}'^{3}$; the $66$ non-zero weights form three
      Galois orbits with $(\text{size},\text{multiplicity})=(6,3),(12,1),(12,3)$.
\item \textbf{Theorem~\ref{thm:master}} (the master formula). For every $S\subseteq C$,
      $\dim\z(S)=12+\sum\{m_\lambda:\lambda|_S=0\}$. The function is the flat-function of $30$
      hyperplanes in a $4$-space.
\item \textbf{Theorem~\ref{thm:spectrum}} (the image). The arrangement has $109$ flats, and the
      image is the eleven-element set $\{12,14,16,18,20,26,28,30,\allowbreak 36,46,78\}$ ---
      every value attained.
\item \textbf{Theorem~\ref{thm:46}} (arithmetic inaccessibility). On the $(x_8,x_{16})$-plane,
      $\dim\z(a x_8+b x_{16})=30$ unless $-b/a$ is a root of an explicit irreducible cubic
      $g\in\Q[t]$, in which case it is $46$. Since $g$ is irreducible over $\Q$, \emph{every}
      rational direction in that plane gives $30$. The multiplicity $16=46-30$ is derived, as
      the exponent in $\charpoly(Q)=c\,g(t)^{16}$, not recorded as a difference.
\item \textbf{Corollary~\ref{cor:sampling}} (why the sample failed). A sample of rational
      directions returns the generic value with probability one on each stratum it does not meet
      by design; the prior sample's $\{12,30,78\}$ is exactly what the coordinate subspaces
      force, and the other eight values are invisible to it. The sample was not unlucky --- it
      was the wrong instrument.
\end{itemize}

\subsection{What is certified how}\label{sec:certification}

We separate two grades of evidence and never merge them.

\begin{itemize}[leftmargin=1.4em]
\item \textbf{Exact over $\Q$:} Proposition~\ref{prop:abelian}; Theorem~\ref{thm:structure} in
      full; Theorem~\ref{thm:master}; Theorem~\ref{thm:46} including the cubic, its
      irreducibility, its discriminant and the field it generates.
\item \textbf{Exhaustive at three faithful primes} ($p=409,421,487$): the count of $109$ flats
      and hence the eleven-element image of Theorem~\ref{thm:spectrum}. This is \emph{not} a
      $\overline{\Q}$ certificate. Reduction modulo $p$ can only \emph{add} linear dependencies
      among weights, so a flat computed mod $p$ could in principle be coarser than its
      characteristic-zero counterpart. Three independent faithful primes agree on all $109$
      flats, and six of the eleven values are corroborated by independent computations sharing
      no code with this one; the exact-over-$\overline{\Q}$ flat lattice is nevertheless
      registered as an open residue rather than claimed. See Scope~\ref{sc:primes}.
\end{itemize}

\section{The structure theorem}\label{sec:structure}

\begin{theorem}[structure of the adjoint action of $C$]\label{thm:structure}
With $C$ as in Definition~\ref{def:C}:
\begin{enumerate}[label=\textup{(\roman*)},leftmargin=2.4em]
\item $\dim\z(C)=12$, and the generalised zero-eigenspace of $\ad(C)$ \emph{equals} $\z(C)$;
      equivalently $C$ acts as literally zero on $\z(C)$, not merely nilpotently.
\item $\esix=\z(C)\oplus V'$ as $C$-modules, with $\dim V'=66$.
\item For generic $R\in\ad(C)$,
      $\charpoly(R)=t^{12}\,q_6(t)^{3}\,q_{12}(t)\,q_{12}'(t)^{3}$, where $q_6$ has degree $6$
      and $q_{12},q_{12}'$ have degree $12$.
\item The $66$ non-zero weights of $C$ on $V'$ consist of $30$ distinct linear functionals on
      $C\cong\Q^4$, falling into three Galois orbits:
      \[
        6\text{ weights of multiplicity }3,\quad
        12\text{ of multiplicity }1,\quad
        12\text{ of multiplicity }3,
      \]
      accounting for $18+12+36=66$ dimensions, and $12+66=78=\dim\esix$.
\end{enumerate}
All four statements are exact over $\Q$.
\end{theorem}

Statement (i) is the one that makes everything else work, and it is not automatic: for a general
abelian subalgebra the generalised zero-eigenspace is strictly larger than the centraliser, and
then the weight decomposition acquires Jordan blocks and the counting below fails. Here the
action on $\z(C)$ is identically zero, so $V'$ is a genuine sum of weight spaces.

\begin{proof}[Method]
The four charges are built explicitly over $\Q$ from $\slt$-lowering applied to highest-weight
vectors, in a fixed Chevalley basis of $\esix$ with $72$ roots. Brackets, the nullspace computing
$\z(C)$, the generalised zero-eigenspace, and the characteristic polynomial of a generic member
of $\ad(C)$ are all computed symbolically over $\Q$; factorisation of the characteristic
polynomial into $t^{12}q_6^3q_{12}q_{12}'^3$ and the resulting orbit structure are exact. Each
step is gated by a control that must be able to fail --- among them, that the six weights of
total multiplicity $18$ vanishing on the $(x_8,x_{16})$-plane reproduce an independently banked
value $12+18=30$ that was not chosen here. See Appendix~\ref{app:verify}.
\end{proof}

\section{The master formula}\label{sec:master}

\begin{theorem}[the rung function is a flat-function]\label{thm:master}
Let $\Lambda\subset C^*$ be the set of $30$ non-zero weights of Theorem~\ref{thm:structure}(iv),
with multiplicities $m_\lambda$. Then for every subspace $S\subseteq C$,
\[
   \dim\z(S)\;=\;12+\sum_{\substack{\lambda\in\Lambda\\ \lambda|_S=0}} m_\lambda .
\]
\end{theorem}

\begin{proof}
Decompose $\esix=\z(C)\oplus\bigoplus_{\lambda\in\Lambda}V_\lambda$ with $\dim
V_\lambda=m_\lambda$; this is a decomposition into $C$-weight spaces by
Theorem~\ref{thm:structure}(i)--(ii), the point of (i) being that no generalised eigenspace is
larger than its eigenspace. An element $X\in V_\lambda$ satisfies $[X,S]=0$ if and only if
$\lambda$ annihilates $S$, and $\z(C)$ is centralised by every $S\subseteq C$. Summing
dimensions gives the formula.
\end{proof}

\begin{corollary}\label{cor:arrangement}
Each $\lambda\in\Lambda$ defines a hyperplane $H_\lambda=\ker\lambda$ in $C\cong\Q^4$. The
value $\dim\z(S)$ depends on $S$ only through the set $\{\lambda:S\subseteq H_\lambda\}$, i.e.\
only through the flat of the arrangement $\{H_\lambda\}_{\lambda\in\Lambda}$ spanned by $S$.
The infinite subspace lattice therefore has finite image under $\dim\z(-)$, and that image is
computed by enumerating flats.
\end{corollary}

This is the paper's structural point. The question ``what values does $\dim\z(S)$ take?'' is a
question about an infinite lattice; Corollary~\ref{cor:arrangement} converts it into a question
about a finite arrangement, which is a finite computation.

\section{The image}\label{sec:image}

\begin{theorem}[the eleven values, all attained]\label{thm:spectrum}
The arrangement of Corollary~\ref{cor:arrangement} has $109$ flats, and
\[
   \bigl\{\dim\z(S)\;:\;S\subseteq C\bigr\}
   =\{12,\,14,\,16,\,18,\,20,\,26,\,28,\,30,\,36,\,46,\,78\},
\]
with each value attained at the subspace dimensions shown:
\begin{center}
\begin{tabular}{@{}lccccccccccc@{}}
\toprule
$\dim\z(S)$      & 12 & 14 & 16 & 18 & 20 & 26 & 28 & 30 & 36 & 46 & 78\\
attained at $\dim S=$ & 4 & 3 & 2 & 2,3 & 1,2 & 2 & 1 & 2 & 1 & 1 & 0\\
\bottomrule
\end{tabular}
\end{center}
The enumeration is exhaustive at each of the three faithful primes $409,421,487$, which agree.
\end{theorem}

\begin{scope}[the grade of this statement]\label{sc:primes}
A prime $p$ is called \emph{faithful} here when the reduction of the weight system mod $p$ still
has $30$ distinct weights with multiplicity profile matching the exact $\Q$-orbits; this is
checked, not assumed, and is a control that can fail. Given faithfulness, enumeration of flats
mod $p$ is exhaustive. What is \emph{not} established is that no flat is coarser than its
characteristic-zero counterpart: reduction can only add dependencies among weights, so the mod-$p$
flat lattice could in principle merge two characteristic-zero flats. Three independent faithful
primes agreeing on all $109$ flats makes this unlikely but does not exclude it. We therefore
state Theorem~\ref{thm:spectrum} at this grade and register the $\overline{\Q}$ certificate as
open. The endpoints are unaffected: $78$ at $S=0$ and $12$ at $S=C$ are exact over $\Q$ by
Theorem~\ref{thm:structure}.
\end{scope}

\begin{remark}[which values are exposed, and in which direction]\label{rem:exposed}
The residue of Scope~\ref{sc:primes} can be stated more precisely than ``not a $\overline{\Q}$
certificate'', and it is worth doing so because the exposure is one-sided.

\textbf{Direction.} Reduction modulo $p$ can only \emph{add} linear dependencies among weights.
More weights vanishing on a subspace means a larger $\sum m_\lambda$, hence a \emph{larger}
$\dim\z(S)$. So a mod-$p$ enumeration can \textbf{contain a value that does not occur over
$\overline{\Q}$; it cannot lose one that does.} The exposed claim is therefore \emph{attainment},
not the bound --- Theorem~\ref{thm:structure}'s containment is exact over $\Q$ and is untouched.

\textbf{Which values.} Six of the eleven --- $12$, $14$, $18$, $30$, $46$ and (trivially) $78$ ---
were obtained in earlier computations sharing no code with the enumeration here, so a mod-$p$
coincidence cannot be manufacturing them. \textbf{The remaining five are exposed:}
\[
   16,\quad 20,\quad 26,\quad 28,\quad 36 .
\]
A $\overline{\Q}$ certificate is needed only for these, and anyone attempting one should check
exactly them.

\textbf{Why such a certificate is not a routine computation.} A value is certified attained by
exhibiting one subspace realising it, and $\dim\z(S)$ is computable exactly over $\Q$ for any
rational $S$ by a single nullspace. We searched for rational witnesses to the five and found
none: no subspace of dimension $1$ or $2$ spanned by vectors with coefficients in
$\{-2,\dots,2\}$, respectively $\{-1,0,1\}$, realises any of $16,20,26,28,36$. This is a bounded
search and therefore evidence rather than proof, but it points the same way as
Theorem~\ref{thm:46}: the exposed values, like the $46$, appear to be reached only at irrational
directions, so a certificate must be assembled in the splitting field of $q_6q_{12}q_{12}'$ rather
than over $\Q$. \textbf{The difficulty of the residue is thus of a piece with the paper's own
thesis} --- these values are hidden from rational access, which is exactly why a rational sample
never found them.
\end{remark}

\section{The $46$, and arithmetic inaccessibility}\label{sec:46}

The value most resistant to sampling is $46$, and its resistance has an exact explanation.

\begin{theorem}[the distinguished plane]\label{thm:46}
Let $W=\im(\ad x_{16})$, so $\dim W=48$; $W$ is invariant under both $\ad x_8$ and $\ad x_{16}$,
and $\ad x_{16}|_W$ is invertible. Put $Q=(\ad x_{16}|_W)^{-1}\circ(\ad x_8|_W)$. Then
\[
   \charpoly(Q)=c\cdot g(t)^{16}
\]
for a constant $c$ and a single cubic $g\in\Q[t]$, and for $(a,b)\ne(0,0)$,
\[
   \dim\z(a x_8+b x_{16})=
   \begin{cases}
      46, & g(-b/a)=0,\\
      30, & \text{otherwise.}
   \end{cases}
\]
Moreover $g$ is irreducible over $\Q$; its discriminant has squarefree part $77$, the same as
that of $x^3-12x-5$; and $g$ acquires a root in $K=\Q[x]/(x^3-12x-5)$, over which it factors
with degrees $(1,2)$. Hence $g$ generates $K$.
\end{theorem}

\begin{corollary}[inaccessible, not merely rare]\label{cor:inaccessible}
Since $g$ is irreducible over $\Q$, it has no rational root, so \emph{every} rational direction
in the $(x_8,x_{16})$-plane gives $\dim\z=30$. The value $46$ is not a rare event for rational
sampling in this plane --- it is an impossible one. Detecting it requires either exact work over
an extension, or the structural route of Theorem~\ref{thm:master}.
\end{corollary}

\begin{remark}[the multiplicity is derived]
The exponent $16$ in $\charpoly(Q)=c\,g^{16}$ is computed, and it equals the jump $46-30$. An
earlier treatment obtained $46$ from a $55$-digit numerical spectral gap and left an independent
determination of the jump as an open item; here the jump is the exponent, and the value is
exact.
\end{remark}

\section{Why the sample returned three values}\label{sec:sampling}

\begin{corollary}[the sample was the wrong instrument]\label{cor:sampling}
A prior investigation evaluated $\dim\z(S)$ on sixteen coordinate subspaces of $C$ together with
$440$ random rational directions, obtaining the image $\{12,30,78\}$, and concluded on that
evidence that the realized spectrum ``appears to be far smaller'' than the eleven-element bound.
Theorem~\ref{thm:spectrum} shows the bound is in fact attained. The discrepancy is fully
explained:
\begin{enumerate}[label=\textup{(\roman*)},leftmargin=2.4em]
\item $78$ is $S=0$ and $12$ is generic, so both are forced;
\item $30$ is the value on the $(x_8,x_{16})$-plane, one of the sixteen coordinate subspaces
      tried;
\item each of the remaining eight values is supported on a proper subvariety of the relevant
      Grassmannian. A random rational direction lies on such a subvariety with probability zero,
      and by Corollary~\ref{cor:inaccessible} the case of $46$ is worse than probability zero:
      it is arithmetically impossible over $\Q$ in the plane where it occurs.
\end{enumerate}
Enlarging the sample does not help, since (iii) is not a matter of sample size.
\end{corollary}

We record this because the failure mode is general and is easy to repeat. Whenever an invariant
is upper semicontinuous on a parameter space, sampling estimates its \emph{generic} value and
provides no information about its image. The remedy used here --- identify the stratification,
then enumerate it --- is available whenever the jumps are governed by a finite arrangement, which
Theorem~\ref{thm:master} establishes in this case.

\begin{remark}[what this settles downstream]
A statement of the form ``\emph{if} a locus with $\dim\z(S)=14$ occurs, its type is forced''
has its antecedent discharged by Theorem~\ref{thm:spectrum}: $14$ is attained, at
$\dim S=3$.
\end{remark}

\section{What is not claimed}

Nothing here concerns real forms, conjugacy classes of subalgebras, or any structure beyond the
dimensions of centralisers of subspaces of the single $C$ of Definition~\ref{def:C}. The
$\overline{\Q}$ flat lattice is open (Scope~\ref{sc:primes}). No interpretation of any of these
dimensions outside Lie theory is offered or used.

\appendix

\section{Verification}\label{app:verify}

The computation is reproduced by \texttt{verify/rung\_attained.py}, which builds $\esix$ in a
Chevalley basis with $72$ roots, constructs $C$ from Definition~\ref{def:C} over $\Q$, and runs
every step behind a control that is designed so that a plausible error makes it fail. The script
exits non-zero on any failure. Selected controls, in the order they run:

\begin{center}
\begin{tabular}{@{}lp{9.2cm}@{}}
\toprule
discharges & control \\
\midrule
\S\ref{sec:C}          & $\esix$ carrier has $72$ roots and dimension $78$; the four charges are built at degrees $8,14,16,22$; all six pairwise brackets vanish \\
Prop.~\ref{prop:forced} & (\texttt{verify/check\_forcing.py}, $13$ checks) $t$ and $\Phi$ are $2T$-invariant, with a control that a non-invariant degree-$8$ form is rejected; $D(2T)\cap E=\{8,14,16,22\}$ exactly; each degree has a unique monomial; controls that $2O$ and $2I$ give \emph{different} intersections, so the McKay partner is load-bearing \\
Prop.~\ref{prop:abelian} & $\dim\z(C)=12$, exactly over $\Q$, reproducing two independent earlier computations \\
Thm~\ref{thm:structure}  & the zero weight has multiplicity $12=\dim\z(C)$; the generalised zero-space \emph{is} $\z(C)$; $C$ acts as literally zero on it; dimensions account exactly to $78$; thirty non-zero weights in three Galois orbits \\
Thm~\ref{thm:46}         & $\dim W=48$; $W$ is $\ad x_8$- and $\ad x_{16}$-invariant; $\ad x_{16}|_W$ invertible; $\charpoly(Q)$ is a single cubic to the $16$th power; that cubic is irreducible over $\Q$; its discriminant squarefree part is $77$; it acquires a root in $K$ with factor degrees $(1,2)$ \\
Thm~\ref{thm:spectrum}   & each prime is checked \emph{faithful} before use ($30$ distinct weights, multiplicity profile $(1,12),(3,18)$ matching the exact orbits); the $(x_8,x_{16})$-plane control returns exactly six weights of total multiplicity $18$, giving $12+18=30$ --- a value fixed independently and not chosen here \\
Cor.~\ref{cor:sampling}  & the sixteen coordinate subspaces reproduce the earlier sample's table $\{12,30,78\}$ exactly --- the sample this paper replaces is regenerated rather than quoted \\
\bottomrule
\end{tabular}
\end{center}

The last two rows are the ones worth attention. The plane control fixes a number the enumeration
did not get to choose, and the sample control reproduces the very evidence whose interpretation
this paper overturns, rather than citing it.

\section{Status in the programme register}\label{app:register}

This is the second of four papers drawn from a register of questions carrying explicit statuses.
The entries this paper bears on:

\begin{center}
\begin{tabular}{@{}llp{7.0cm}@{}}
\toprule
entry & status & relation to this paper \\
\midrule
\texttt{OA-C0005} & \texttt{PROVED} & that the binary tetrahedral data determines the $E_6$ root-system type. Used as background; not re-derived here. \\
\texttt{OA-C0006} & \texttt{CONDITIONAL} & whether the principal $\slt$ placement is selected. \textbf{Assumed, not established.} $C$ is \emph{defined} using the principal $\slt$; this paper does not argue that placement is forced. \\
\texttt{OA-C1057} & \texttt{PROVED} & exact triples for distinguished non-regular $E_6$ characteristics; consistent with, and independent of, the enumeration here. \\
\bottomrule
\end{tabular}
\end{center}

\noindent
The honest reading is that this paper's theorems are unconditional statements about a
\emph{given} $C$, while the interest of that particular $C$ rests on \texttt{OA-C0006}, which is
\texttt{CONDITIONAL}. We state the theorems in a form that does not depend on resolving it.

\begin{thebibliography}{9}
\bibitem{Bourbaki} N.~Bourbaki, \emph{Groupes et alg\`ebres de Lie}, Ch.~IV--VI, Hermann, 1968.
\bibitem{Carter} R.~W.~Carter, \emph{Finite Groups of Lie Type: Conjugacy Classes and Complex
  Characters}, Wiley, 1985.
\bibitem{CollingwoodMcGovern} D.~H.~Collingwood and W.~M.~McGovern, \emph{Nilpotent Orbits in
  Semisimple Lie Algebras}, Van Nostrand Reinhold, 1993.
\bibitem{McKay} J.~McKay, Graphs, singularities, and finite groups, in \emph{The Santa Cruz
  Conference on Finite Groups}, Proc.\ Sympos.\ Pure Math.\ \textbf{37}, AMS, 1980, 183--186.
\bibitem{Klein} F.~Klein, \emph{Vorlesungen \"uber das Ikosaeder und die Aufl\"osung der
  Gleichungen vom f\"unften Grade}, Teubner, 1884.
\bibitem{OrlikTerao} P.~Orlik and H.~Terao, \emph{Arrangements of Hyperplanes}, Springer, 1992.
\bibitem{Kostant} B.~Kostant, The principal three-dimensional subgroup and the Betti numbers of
  a complex simple Lie group, \emph{Amer.\ J.\ Math.} \textbf{81} (1959), 973--1032.
\end{thebibliography}

\end{document}
